\begin{lemma}\label{hessian:sobolev-codimension-two}
Let $E$ be a compact subset of a smooth real submanifold of
codimension two contained in a coordinate chart. Then
$\operatorname{Cap}_1(E)=0$.
\end{lemma}
\begin{proof}
It suffices to treat a compact subset of a submanifold chart, since
finitely many such subsets cover $E$ and the capacity is subadditive.
Choose real coordinates $(s,t)$ there, with $t\in\mathbb R^2$ and the
submanifold given by $t=0$. Choose $0\leq\eta\leq1$ smooth,
supported in this chart and equal to one near the compact subset.
Put $r=|t|$.
For $0<\epsilon<1/2$ define
\[
 h_\epsilon(r)=
 \begin{cases}
  1,&0\leq r\leq\epsilon^2,\\
  \displaystyle\frac{\log(\epsilon/r)}{\log(1/\epsilon)},
       &\epsilon^2<r<\epsilon,\\
  0,&r\geq\epsilon.
 \end{cases}
\]
The Sobolev function $f_\epsilon=\eta h_\epsilon$ equals one on an
open neighborhood of that subset. On the fixed compact coordinate
set the Sobolev norm is equivalent to the Euclidean one. Integration
in the two normal variables gives
\[
 \|f_\epsilon\|_{H^1}^2
 \leq C\left(\epsilon^2+
     \int_{\epsilon^2}^{\epsilon}
       \frac{\mathrm dr}{r\log(1/\epsilon)^2}\right)
 =C\left(\epsilon^2+\frac1{\log(1/\epsilon)}\right)
 \longrightarrow0.
\]
The derivatives of $\eta$ contribute at most $C\epsilon^2$.
Taking the infimum over open neighborhoods proves the assertion.
\end{proof}

\begin{example}\label{hessian:sobolev-two-product-obstruction}
There are bounded continuous $\omega$-plurisubharmonic potentials
$u,v$ and a compact Borel set $E$ such that
\[
 \operatorname{Cap}_1(E)=0,\qquad
 \int_E Q_2(u,v)\wedge\omega^{m-2}>0.
\]
Consequently the first-Hessian capacity estimate does not extend to
this positive two-factor measure, even within the continuous
$\omega$-plurisubharmonic subclass.
\end{example}
\begin{proof}
Choose holomorphic coordinates $z_j=x_j+\mathrm i y_j$ and a smooth
nonnegative cutoff $\chi$ supported in the chart and equal to one
on a smaller coordinate box. Extend
\[
 u=a\chi|x_1|,\qquad v=a\chi|x_2|
\]
by zero outside the chart. These functions are globally continuous
and bounded. To check their curvature, replace $|x_j|$ by
$g_\delta(x_j)=(x_j^2+\delta^2)^{1/2}$, with $0<\delta\leq1$.
The Hessian of $\chi g_\delta$ has the nonnegative term
$\chi\ddc g_\delta$; all its remaining terms are uniformly bounded
forms, since $g_\delta$ and $g_\delta'$ are uniformly bounded on the
fixed support. Thus $\ddc(\chi g_\delta)\geq-C\omega$ with $C$
independent of $\delta$. Choosing $a>0$ sufficiently small gives
$\omega+\ddc(a\chi g_\delta)\geq\omega/2$ for both coordinates.
Their decreasing limits $u,v$ are therefore
$\omega$-plurisubharmonic and admissible.

On the box where $\chi=1$, set
$\alpha_j=\frac{\mathrm i}{4\pi}\mathrm dz_j\wedge\mathrm d\bar z_j$.
The book's normalization gives
\[
 \ddc g_\delta(x_j)=\tfrac12g_\delta''(x_j)\alpha_j,
 \qquad g_\delta''(x_j)
       =\frac{\delta^2}{(x_j^2+\delta^2)^{3/2}}
       \longrightarrow2\delta_{\{x_j=0\}}.
\]
The measures $g_\delta^{\prime\prime}(t)\,\mathrm dt$ have mass two,
while their mass outside any fixed neighborhood of zero tends to
zero. The two
smooth factors depend on distinct real variables, so their product
is an approximate identity in $(x_1,x_2)$ and converges to the
tensor-product measure. By
\cref{hessian:bounded-two-factor-convergence}, locally
\[
 D_2(u,v)=a^2\delta_{\{x_1=0\}}\delta_{\{x_2=0\}}
                         \alpha_1\wedge\alpha_2.
\]
Choose $E$ to be a compact piece of $\{x_1=x_2=0\}$ inside this
box with nonempty relative interior. The preceding lemma gives
$\operatorname{Cap}_1(E)=0$. The positive measure
$D_2(u,v)\wedge\omega^{n-2}$ has strictly positive mass on $E$,
since $\omega$ is strictly positive. The smooth term $\omega^n$ and
the two single-Hessian terms in the expansion
\[
 Q_2(u,v)\wedge\omega^{m-2}
 =\omega^n+\omega^{n-1}\wedge(\ddc u+\ddc v)
                         +D_2(u,v)\wedge\omega^{n-2}
\]
have zero mass on $E$: each single-Hessian term is absolutely
continuous in the other normal real variable. Hence the asserted
mass is strictly positive. A bound by
$C\operatorname{Cap}_1(F)^b$ for all Borel $F$, with any $b>0$,
would contradict this choice of $E$.
\end{proof}
